% GENERATED FILE. Edit canonical lessons, then run npm run generate:books.
% canonical-source-hash: fnv1a64:89579660ce718429
% canonical-lessons: ZH-C285-zhaohui, ZH-C285-problem, ZH-C285-zhengjian, ZH-C285-guashi, ZH-C285-lianxi, ZH-C285-potential, ZH-W285-a-note, ZH-C285-practice

\chapter{Lost on the Way}
\label{ch:zh-lost-on-the-way}
\hlchaptermodality{285}

\begin{chapteropening}
\textbf{By the end of this chapter:} \emph{I can report a lost thing and write a short note about it.}

Reporting a problem with …\zh{不见了}, saying you can't find it with \zh{找不到}, papers, reporting a loss and getting in touch, and writing four connected lines to a hotel desk.
\end{chapteropening}

\section[\texorpdfstring{\zh{找回} (zhǎohuí)}{找回 (zhǎohuí)}]{\zh{找回} (zhǎohuí) — to get back}

\label{lesson:ZH-C285-zhaohui}



\begin{quote}
A new chapter, and it opens on the happy end of a bad day.

\begin{itemize}
\raggedright
  \item \emph{Recall:} say \zh{我们坐八点的火车} — \emph{wǒmen zuò bā diǎn de huǒchē} — \textbf{R1}, the chapter you have finished
  \item \emph{Recall:} say \zh{找} — \emph{zhǎo}, to look for — \textbf{R3}, durable
  \item \emph{Recall:} say \zh{找到} — \emph{zhǎodào}, to find — \textbf{R3}
  \item \emph{Recall:} say \zh{回} — \emph{huí}, to go back — \textbf{R3}
\end{itemize}
\end{quote}

\subsection*{You'll want to know: \zh{找回}}
\textbf{\zh{找回}} — \emph{zhǎohuí} — \textquotedblleft{}to get back\textquotedblright{}, something you had lost.

\begin{quote}
\textbf{\zh{我找回了我的手机}} \emph{wǒ zhǎohuí le wǒ de shǒujī} — I got my phone back.
\end{quote}

Both characters are ones you write.

\begin{cousinweb}
\textbf{\zh{找回}} (\emph{zhǎohuí}) is built from \zh{找} 'to look for' and \zh{回} 'to go back'. The second is inherited from Middle Chinese \emph{hwoj} (Guangyun, in Baxter's transcription), glossed 'to turn round'; the dictionary has no entry for \zh{找}, so the history of that part is not settled.
\end{cousinweb}

\begin{grammarlens}[title={What you've built}]
\zh{找} is to look for, \zh{找到} is to find, and \zh{找回} is to find what you had lost and have it again.

\zh{了} after it says it is done: \zh{我找回了我的手机}.
\end{grammarlens}

\subsection*{Your turn}
\begin{itemize}
\raggedright
  \item \emph{Say it:} \emph{zhǎohuí}
  \item \emph{Say it:} \emph{zhǎohuí}, once more
  \item \emph{Say it:} \zh{我找回了我的手机} — \emph{wǒ zhǎohuí le wǒ de shǒujī} — I got my phone back.
  \item \emph{Recall:} say \zh{找}, then \zh{找到}, then \zh{找回}
\end{itemize}

\begin{tcolorbox}[breakable,colback=teal!4,colframe=teal!35!black,title={Before you move on}]
What does \zh{找回} mean? (\textquotedblleft{}To get back\textquotedblright{}.) And \zh{找到}? (\textbf{To find}.)
\end{tcolorbox}
\section[\texorpdfstring{\zh{我的手机不见了} (wǒ de shǒujī bú jiàn le)}{我的手机不见了 (wǒ de shǒujī bú jiàn le)}]{\zh{我的手机不见了} — it's gone}

\label{lesson:ZH-C285-problem}



\begin{quote}
Every word in it is one you have; together they report a problem.

\begin{itemize}
\raggedright
  \item \emph{Read it:} \zh{找回}, aloud — \emph{zhǎohuí}
  \item \emph{Recall:} say \emph{diū}, to lose — \textbf{R3}, durable
  \item \emph{Recall:} say \zh{手机} — \emph{shǒujī}, a phone — \textbf{R3}
\end{itemize}
\end{quote}

\begin{grammarlens}[title={Grammar Lens: …\zh{不见了}, it has gone missing}]
\begin{quote}
\textbf{Name the thing, then say \zh{不见了}: it has gone missing.}
\end{quote}

\begin{quote}
\textbf{\zh{我的手机不见了}} \emph{wǒ de shǒujī bú jiàn le} — My phone's gone missing. \textbf{\zh{我的车票不见了} \zh{你看到了吗}} \emph{wǒ de chēpiào bú jiàn le, nǐ kàndào le ma} — My ticket's disappeared. Have you seen it?
\end{quote}

\zh{不见} is 'not seen', and \zh{了} says that is new: it was there, and now it is not. \zh{不} before the fall of jiàn is \emph{bú}.

\emph{diū le} says you lost it yourself: \emph{wǒ de qiánbāo diū le}. \emph{huài le} says it broke. \emph{qǐng bāng wǒ}, please help me, asks for help.
\end{grammarlens}

\subsection*{Your turn}
Say these aloud:

\begin{itemize}
\raggedright
  \item \zh{我的手机不见了} — \emph{wǒ de shǒujī bú jiàn le} — My phone's gone missing.
  \item \zh{我的车票不见了} \zh{你看到了吗} — \emph{wǒ de chēpiào bú jiàn le, nǐ kàndào le ma} — My ticket's disappeared. Have you seen it?
  \item \emph{qǐng bāng wǒ} — Please help me.
  \item something of yours that went missing once, with \zh{不见了}
\end{itemize}

\begin{tcolorbox}[breakable,colback=teal!4,colframe=teal!35!black,title={Before you move on}]
What comes after the thing? (\textbf{\zh{不见了}}.) How is \zh{不} said there? (\emph{bú}, before the fall of jiàn.)
\end{tcolorbox}
\section[zhèngjiàn]{zhèngjiàn — papers; an identity document}

\label{lesson:ZH-C285-zhengjian}



\begin{quote}
Before the new one: last lesson's problem, and two things you hope never to lose.

\begin{itemize}
\raggedright
  \item \emph{Read it:} \zh{我的手机不见了}, aloud — \emph{wǒ de shǒujī bú jiàn le}
  \item \emph{Recall:} say \emph{hùzhào}, a passport — \textbf{R3}, durable
  \item \emph{Recall:} say \emph{shēnfènzhèng}, an identity card — \textbf{R3}
\end{itemize}
\end{quote}

\subsection*{You'll want to know: zhèngjiàn}
\emph{zhèngjiàn} — \textquotedblleft{}papers; an identity document\textquotedblright{}.

\begin{quote}
\emph{wǒ de zhèngjiàn bú jiàn le} — My papers have gone missing.
\end{quote}

Its characters come later; for now it is a word for the ear.

\begin{cousinweb}
\emph{zhèngjiàn} is built from zhèng 'proof' and jiàn 'an item', both inherited from Middle Chinese, \emph{tsyingH} and \emph{gjenX} (Guangyun, in Baxter's transcription), which glosses the first, under its traditional form, 'to verify' and the second 'a portion'.
\end{cousinweb}

\begin{grammarlens}[title={What you've built}]
\emph{hùzhào} and \emph{shēnfènzhèng} are both \emph{zhèngjiàn}: it is the word for any of them, the one a desk asks to see.
\end{grammarlens}

\subsection*{Your turn}
\begin{itemize}
\raggedright
  \item \emph{Say it:} \emph{zhèngjiàn}
  \item \emph{Say it:} \emph{zhèngjiàn}, once more
  \item \emph{Say it:} \emph{wǒ de zhèngjiàn bú jiàn le} — My papers have gone missing.
  \item \emph{Recall:} say \emph{hùzhào}, then \emph{zhèngjiàn}: one of them, then any of them
\end{itemize}

\begin{tcolorbox}[breakable,colback=teal!4,colframe=teal!35!black,title={Before you move on}]
What does \emph{zhèngjiàn} mean? (\textquotedblleft{}Papers; an identity document\textquotedblright{}.) Is a passport one? (\textbf{Yes}.)
\end{tcolorbox}
\section[guàshī]{guàshī — to report a loss}

\label{lesson:ZH-C285-guashi}



\begin{quote}
Before the new one: last lesson's word, and two words that share its sounds.

\begin{itemize}
\raggedright
  \item \emph{Recall:} say \emph{zhèngjiàn}, papers — \textbf{R1}, one lesson back
  \item \emph{Recall:} say \emph{guà}, to hang — \textbf{R3}, durable
  \item \emph{Recall:} say \emph{shīwàng}, disappointed — \textbf{R3}
\end{itemize}
\end{quote}

\subsection*{You'll want to know: guàshī}
\emph{guàshī} — \textquotedblleft{}to report the loss of\textquotedblright{} a card or a document.

\begin{quote}
\emph{wǒ yào guàshī wǒ de xìnyòngkǎ} — I need to report my credit card lost.
\end{quote}

\begin{quote}
\emph{wǒ děi qù guàshī} — I have to go and report it lost.
\end{quote}

Its characters come later; for now it is a word for the ear.

\begin{cousinweb}
\emph{guàshī} is built from guà 'to hang' and shī 'to lose', both inherited from Middle Chinese, \emph{kweaH} and \emph{syit} (Guangyun, in Baxter's transcription), which glosses the first 'to hang' and the second 'to let slip'. 'To report a loss' for the pair is its own sense, which neither gloss gives.
\end{cousinweb}

\begin{grammarlens}[title={What you've built}]
\emph{guàshī} is what you do at a bank or a police desk once a card or a passport is gone, so nobody else can use it.

The thing follows it: \emph{guàshī wǒ de xìnyòngkǎ}.
\end{grammarlens}

\subsection*{Your turn}
\begin{itemize}
\raggedright
  \item \emph{Say it:} \emph{guàshī}
  \item \emph{Say it:} \emph{guàshī}, once more
  \item \emph{Say it:} \emph{wǒ yào guàshī wǒ de xìnyòngkǎ} — I need to report my credit card lost.
  \item \emph{Say it:} \emph{wǒ děi qù guàshī} — I have to go and report it lost.
  \item \emph{Recall:} say \emph{guà}, then \emph{guàshī}
\end{itemize}

\begin{tcolorbox}[breakable,colback=teal!4,colframe=teal!35!black,title={Before you move on}]
What does \emph{guàshī} mean? (\textquotedblleft{}To report a loss\textquotedblright{}.) What follows it? (\textbf{The thing that was lost}.)
\end{tcolorbox}
\section[liánxì]{liánxì — to get in touch}

\label{lesson:ZH-C285-lianxi}



\begin{quote}
Before the new one: last lesson's word, and how you reach people now.

\begin{itemize}
\raggedright
  \item \emph{Recall:} say \emph{guàshī}, to report a loss — \textbf{R1}, one lesson back
  \item \emph{Recall:} say \zh{打电话} — \emph{dǎ diànhuà}, to phone — \textbf{R3}, durable
  \item \emph{Recall:} say \emph{hàomǎ}, a number — \textbf{R3}
\end{itemize}
\end{quote}

\subsection*{You'll want to know: liánxì}
\emph{liánxì} — \textquotedblleft{}to get in touch with; to contact\textquotedblright{}.

\begin{quote}
\emph{wǒ zěnme liánxì nǐ} — How can I get in touch with you?
\end{quote}

\begin{quote}
\emph{qǐng gēn wǒ liánxì} — Please get in touch with me.
\end{quote}

Its characters come later; for now it is a word for the ear.

\begin{cousinweb}
\emph{liánxì} is built from lián 'to join' and xì 'to tie'. The Guangyun has no entry for the first under the traditional forms we look it up by, so the history of that part is not settled; the second is inherited from Middle Chinese \emph{kejH} (Guangyun, in Baxter's transcription), glossed 'joined together' under one traditional form and 'to tie up' under another.
\end{cousinweb}

\begin{grammarlens}[title={What you've built}]
\zh{打电话} is one way to \emph{liánxì}; a message is another. The person you reach follows it, or comes before it with gēn: \emph{qǐng gēn wǒ liánxì}.
\end{grammarlens}

\subsection*{Your turn}
\begin{itemize}
\raggedright
  \item \emph{Say it:} \emph{liánxì}
  \item \emph{Say it:} \emph{liánxì}, once more
  \item \emph{Say it:} \emph{wǒ zěnme liánxì nǐ} — How can I get in touch with you?
  \item \emph{Say it:} \emph{qǐng gēn wǒ liánxì} — Please get in touch with me.
  \item \emph{Recall:} say \zh{打电话}, then \emph{liánxì}: one way, then any way
\end{itemize}

\begin{tcolorbox}[breakable,colback=teal!4,colframe=teal!35!black,title={Before you move on}]
What does \emph{liánxì} mean? (\textquotedblleft{}To get in touch\textquotedblright{}.) How do you ask how to reach someone? (\emph{wǒ zěnme liánxì nǐ}.)
\end{tcolorbox}
\section[\texorpdfstring{\zh{找不到} (zhǎo bu dào)}{找不到 (zhǎo bu dào)}]{\zh{找不到} — can't find it}

\label{lesson:ZH-C285-potential}



\begin{quote}
Verbs with a result after them, which you have; and a way to say the result is out of reach.

\begin{itemize}
\raggedright
  \item \emph{Recall:} say \emph{liánxì}, to get in touch — \textbf{R1}, one lesson back
  \item \emph{Recall:} say \zh{找到} — \emph{zhǎodào}, to find — \textbf{R3}, durable
  \item \emph{Recall:} say \emph{tīngdǒng}, to understand what you hear — \textbf{R3}
  \item \emph{Recall:} say \zh{看到} — \emph{kàndào}, to see — \textbf{R3}
\end{itemize}
\end{quote}

\begin{grammarlens}[title={Grammar Lens: \zh{找不到} and \zh{找得到}, can't and can}]
\begin{quote}
\textbf{\zh{不} between a verb and its result says you can't manage it; \zh{得} there says you can.}
\end{quote}

\begin{quote}
\textbf{\zh{我找不到我的手机}} \emph{wǒ zhǎo bu dào wǒ de shǒujī} — I can't find my phone. \emph{wǒ tīng bu dǒng} — I can't understand what's being said.
\end{quote}

\zh{找到}, \zh{看到}, \emph{tīngdǒng}: you know verbs with a result after them. Put \zh{不} in the middle, said light, and the result is out of reach: \zh{找不到}. Put \zh{得} there, also light, and it is within reach: \textbf{\zh{你找得到吗}} \emph{nǐ zhǎo de dào ma}?

Keep to the results \emph{dào}, \emph{dǒng} and \emph{wán} for now: \emph{zuò bu wán}, more than you can finish.
\end{grammarlens}

\subsection*{Your turn}
Say these aloud:

\begin{itemize}
\raggedright
  \item \zh{我找不到我的手机} — \emph{wǒ zhǎo bu dào wǒ de shǒujī} — I can't find my phone.
  \item \emph{wǒ tīng bu dǒng} — I can't understand what's being said.
  \item \zh{你找得到吗} — \emph{nǐ zhǎo de dào ma} — Can you find it?
  \item something you can't find today, with \zh{找不到}
\end{itemize}

\begin{tcolorbox}[breakable,colback=teal!4,colframe=teal!35!black,title={Before you move on}]
What does \zh{不} do in \zh{找不到}? (\textbf{Puts the result out of reach: can't find}.) And \zh{得} in \zh{找得到}? (\textbf{Puts it within reach: can find}.)
\end{tcolorbox}
\section[(a note)]{A note to the desk, four lines that lean on each other}

\label{lesson:ZH-W285-a-note}



\begin{quote}
You are staying at a hotel, and something of yours has gone missing. The front desk, \emph{qiántái}, can help, if you tell them clearly.

This lesson writes to them.

\begin{itemize}
\raggedright
  \item \emph{Recall:} say \emph{liánxì}, to get in touch — \textbf{R2}
  \item \emph{Recall:} say \emph{qiántái}, the front desk — \textbf{R3}, durable
\end{itemize}
\end{quote}

\subsection*{What to know first}
A note is connected when each line leans on the one before. Four short lines can hang together; four long ones can fall apart.

Three things do the leaning, and you own all three. \textbf{The time words lead}: \zh{早上}, \zh{回来的时候}, \zh{明天}. \textbf{The subject is said once}: \zh{我} in the first line, and the lines after it start on the time or the verb. \textbf{Each line answers the question the last one raised}: gone, when, what now, how to reach you.

Your reader is the desk. Your purpose is to get the thing back.

\subsection*{Writing — connected composition}
Write \textbf{four lines} to the front desk. No name at the top: begin as you would at the desk.

\begin{itemize}
\raggedright
  \item Line 1 says what is gone, with \zh{不见了}.
  \item Line 2 says when you last had it, with \zh{的时候}.
  \item Line 3 says what you need, with \zh{请}, or what you think you should do, with \zh{我应该}.
  \item Line 4 says how to reach you: \emph{liánxì}, and a time or a room number.
\end{itemize}

A word whose characters you have not written goes in pinyin, inside the line, with its tone marks: \emph{liánxì}, \emph{qiántái}, \emph{hùzhào}. Every character comes from the ones you have learned, each in the order you learned it, each in the same size of square.

No Chinese commas or full stops: leave a space between clauses, as this book prints them.

Write the note. Then read it once as the desk would, top to bottom: could they act on it?

\begin{tcolorbox}[breakable,colback=teal!4,colframe=teal!35!black,title={Before you move on}]
Read your second line on its own. Does it stand up without the first? (\textbf{It should not quite: it says when, and the first says what}.)

Count the lines that start with \zh{我}. If every one does, you wrote a list, not a note.
\end{tcolorbox}
\section[Practice]{\zh{我的车票不见了}}

\label{lesson:ZH-C285-practice}



\begin{quote}
The chapter closes here, on a train, with two things missing.

\begin{itemize}
\raggedright
  \item \emph{Recall:} say \emph{guàshī}, to report a loss — \textbf{R2}
  \item \emph{Recall:} say \emph{zhèngjiàn}, papers — \textbf{R2}
  \item \emph{Recall:} say \zh{我找回了我的手机} — \emph{wǒ zhǎohuí le wǒ de shǒujī} — \textbf{R2}
\end{itemize}
\end{quote}

\subsection*{The exchange}
\begin{quote}
\textbf{A:} \zh{我的车票不见了}
\end{quote}

\begin{quote}
\emph{wǒ de chēpiào bú jiàn le} — My ticket's gone.
\end{quote}

\begin{quote}
\textbf{B:} \zh{你找了吗}
\end{quote}

\begin{quote}
\emph{nǐ zhǎo le ma} — Have you looked for it?
\end{quote}

\begin{quote}
\textbf{A:} \zh{找了} \zh{我找不到}
\end{quote}

\begin{quote}
\emph{zhǎo le, wǒ zhǎo bu dào} — I have. I can't find it.
\end{quote}

\begin{quote}
\textbf{B:} \zh{你上车的时候有吗}
\end{quote}

\begin{quote}
\emph{nǐ shàngchē de shíhou yǒu ma} — Did you have it when you got on?
\end{quote}

\begin{quote}
\textbf{A:} \zh{有} \zh{我坐在这儿的时候还有}
\end{quote}

\begin{quote}
\emph{yǒu, wǒ zuò zài zhèr de shíhou hái yǒu} — Yes. I still had it when I sat down here.
\end{quote}

\begin{quote}
\textbf{B:} \emph{nǐ yīnggāi qù qiántái} — You should go to the desk.
\end{quote}

\begin{quote}
\textbf{A:} \emph{wǒ de zhèngjiàn yě bú jiàn le} — My papers are gone too.
\end{quote}

\begin{quote}
\textbf{B:} \emph{nà nǐ děi mǎshàng guàshī} — Then you'll have to report them lost straight away.
\end{quote}

A answers with the verb alone, \zh{找了}, and B leads with the moment, \zh{你上车的时候}. \zh{还有} says A still had it then.

\begin{tcolorbox}[breakable,skin=enhanced,colback=orange!6,colframe=orange!45!black,boxrule=0.5pt,arc=1mm,left=6pt,right=6pt,top=4pt,bottom=4pt,fonttitle=\bfseries,title={Read this}]
\emph{Read it:} the exchange through once without stopping, then answer the question below
\end{tcolorbox}

\begin{grammarlens}[title={What you've built — when something goes missing}]
What to say, from the bad news to the end of it:

\begin{itemize}
\raggedright
  \item \textbf{\zh{我的车票不见了}} \emph{wǒ de chēpiào bú jiàn le} — it's gone.
  \item \textbf{\zh{我找不到}} \emph{wǒ zhǎo bu dào} — I can't find it.
  \item \emph{guàshī}, \emph{liánxì} — report it, and leave a way to reach you.
  \item \textbf{\zh{我找回了我的手机}} \emph{wǒ zhǎohuí le wǒ de shǒujī} — and, with luck, it comes back.
\end{itemize}

\textbf{Listen.} \emph{Nǐ de shénme bú jiàn le?}

\textbf{Speak.} Report something missing at a desk, in four sentences: what, when, what you need, how to reach you.

\emph{Write it:} \zh{不见了} and \zh{找不到}
\end{grammarlens}

\subsection*{Your turn}
\begin{itemize}
\raggedright
  \item \emph{Say it:} \zh{我的车票不见了} \zh{你看到了吗} — \emph{wǒ de chēpiào bú jiàn le, nǐ kàndào le ma} — My ticket's disappeared. Have you seen it?
  \item \emph{Say it:} \zh{我找不到我的手机} — \emph{wǒ zhǎo bu dào wǒ de shǒujī} — I can't find my phone.
  \item \emph{Say it:} the whole exchange, both parts
  \item \emph{Write it:} \zh{我的手机不见了}
\end{itemize}

\begin{tcolorbox}[breakable,colback=teal!4,colframe=teal!35!black,title={Before you move on}]
Say your ticket has gone. (\textbf{\zh{我的车票不见了}}.) Say you can't find it. (\textbf{\zh{我找不到}}.) Say you got your phone back. (\textbf{\zh{我找回了我的手机}}.)
\end{tcolorbox}
